\chapter{Builders and the runtime tree}
\label{ch:builders}

\begin{aside}
The DSL is the path most users take to construct UI. There is a
second path that does not appear in the chapters so far: the
\emph{builders}. A builder is a small fluent API that produces an
\code{Element} programmatically, without the compile-time
machinery. Builders are what widgets use internally, what
runtime-driven UIs use, and what you reach for when the DSL's
type-state machine is in the way.
\end{aside}

\section{Why builders exist}

The DSL is excellent when the structure is known at compile
time. When the structure depends on data, the DSL fights you:

\begin{itemize}[leftmargin=*]
  \item A list of items whose length depends on a vector.
  \item A box whose border style is chosen by a config setting.
  \item A widget whose internal layout depends on its current
    state.
\end{itemize}

You can do these things in the DSL with \code{when}, the runtime
spellings of pipes, and ranges --- but it gets noisy. A builder
sidesteps the compile-time pipeline entirely:

\begin{lstlisting}
auto e = box()
    .direction(Direction::Column)
    .border(BorderStyle::Single)
    .padding(1)
    .add(text("Title").bold())
    .add(text("Body line"))
    .build();
\end{lstlisting}

Each method returns a reference to the builder, so calls chain.
\code{build()} produces the final \code{Element}.

\subsection*{Builders are not anti-DSL}

The two work together. A widget can have its own builder
internally and the DSL call site is unchanged:

\begin{lstlisting}
auto bar = widget(StatusBar()
    .mode(\"NORMAL\")
    .file(\"~/notes.md\")
    .cursor(line, col));
\end{lstlisting}

Here \code{StatusBar} is built with a builder, but the result is
hosted in the DSL by \code{widget(\ldots)}.

\section{The BoxBuilder}

\file{maya/include/maya/element/builder.hpp} defines
\code{BoxBuilder}:

\begin{lstlisting}
class BoxBuilder {
public:
    BoxBuilder& direction(Direction);
    BoxBuilder& justify(Justify);
    BoxBuilder& align(Align);
    BoxBuilder& wrap(Wrap);
    BoxBuilder& gap(int);
    BoxBuilder& padding(Padding);
    BoxBuilder& padding(int);             // all sides
    BoxBuilder& padding(int v, int h);    // vert + horiz
    BoxBuilder& border(BorderStyle);
    BoxBuilder& border_color(Color);
    BoxBuilder& width(int);
    BoxBuilder& height(int);
    BoxBuilder& grow(int);
    BoxBuilder& shrink(int);
    BoxBuilder& style(Style);
    BoxBuilder& add(Element child);
    template <class... Cs> BoxBuilder& add_all(Cs&&... cs);

    Element build() &&;
};

BoxBuilder box();
\end{lstlisting}

Each setter returns a reference to itself, so calls compose.
\code{build()} is rvalue-qualified --- you can only call it on a
temporary (i.e. at the end of the chain), which discourages
accidentally reusing a half-built builder.

\subsection*{Why ampersand-ampersand on build()?}

The double-ampersand (\code{\&\&}) qualifier means ``this method
can only be called on an rvalue''. \code{build()} consumes the
builder. If you wrote:

\begin{lstlisting}
BoxBuilder b;
b.direction(Direction::Row);
auto e = b.build();   // error: build() is not callable on lvalues
\end{lstlisting}

\ldots{} the compiler would refuse. The fix:

\begin{lstlisting}
auto e = std::move(b).build();
\end{lstlisting}

This is a deliberate design: a builder is a one-shot value. The
\code{\&\&} qualifier enforces that.

\begin{funfact}
Ref-qualifiers (\code{\&} and \code{\&\&} on methods) were added
in C++11 specifically for builder patterns and string-builder
optimisations. \code{std::string::operator+=} can be optimised
when it knows it is being called on a temporary; the same idea
applies to fluent builders.
\end{funfact}

\section{The TextBuilder}

For text leaves:

\begin{lstlisting}
class TextBuilder {
public:
    TextBuilder& style(Style);
    TextBuilder& fg(Color);
    TextBuilder& bg(Color);
    TextBuilder& bold();
    TextBuilder& italic();
    TextBuilder& underline();
    TextBuilder& dim();
    TextBuilder& hyperlink(std::string);

    Element build() &&;
    operator Element() && { return std::move(*this).build(); }
};

TextBuilder text(std::string s);
TextBuilder text(std::string_view sv);
\end{lstlisting}

Most users never call \code{build()} on a \code{TextBuilder}
directly; the conversion operator does it implicitly when the
builder is consumed by another builder's \code{add()}:

\begin{lstlisting}
auto e = box()
    .add(text(\"hi\").bold())     // TextBuilder -> Element
    .build();
\end{lstlisting}

The \code{operator Element()} is rvalue-qualified for the same
reason as \code{build()}: it consumes the builder.

\section{ListBuilder for dynamic children}

When you have a runtime list of items:

\begin{lstlisting}
auto e = list()
    .for_each(items, [](auto& item) {
        return text(item.label).bold();
    })
    .with_separator(text(\" | \").dim())
    .build();
\end{lstlisting}

\code{ListBuilder::for\_each} takes a range and a function from
item to element. It eats the runtime cost (one allocation for
the children vector) so the DSL does not have to.

\subsection*{Implementation sketch}

\begin{lstlisting}
template <class R, class F>
ListBuilder& ListBuilder::for_each(R&& range, F&& f) {
    for (auto& item : range) {
        children_.push_back(Element{std::forward<F>(f)(item)});
    }
    return *this;
}
\end{lstlisting}

Generic over the range type and the projection. The result is
sequential children in the order of iteration.

\section{Widget builders}

Every widget in \file{maya/include/maya/widget/} is a builder
whose \code{build()} returns an \code{Element}. The pattern:

\begin{lstlisting}
class Progress {
public:
    Progress& value(double v);                 // 0..1
    Progress& total(int t);
    Progress& current(int c);
    Progress& label(std::string);
    Progress& width(int);
    Progress& style(Style);

    Element build() &&;
};
\end{lstlisting}

The widget knows how to render itself; the caller picks the
parameters. The DSL hosts the result:

\begin{lstlisting}
auto bar = widget(Progress{}
    .total(100)
    .current(progress_now)
    .label(\"Loading\\ldots\")
    .width(40));
\end{lstlisting}

\subsection*{Why fluent, not designated init?}

You can also write widgets with designated initialisers:

\begin{lstlisting}
auto bar = widget(Progress{
    .total = 100,
    .current = progress_now,
    .label = \"Loading\\ldots\",
    .width = 40,
});
\end{lstlisting}

Both work. The fluent form composes nicely with conditional
chains:

\begin{lstlisting}
auto bar = Progress{}.total(100).current(p);
if (show_label) bar.label(\"Loading\");
auto e = widget(std::move(bar));
\end{lstlisting}

The designated-init form requires you to commit to all fields at
once. Pick by ergonomics.

\section{Builders for tree construction}

A common pattern: build a tree by recursion.

\begin{lstlisting}
Element render_node(const Node& n) {
    auto b = box()
        .direction(Direction::Column)
        .add(text(n.label).bold());
    for (auto& child : n.children) {
        b.add(render_node(child));
    }
    return std::move(b).build();
}
\end{lstlisting}

Each call to \code{render\_node} produces an \code{Element}. The
recursion handles arbitrarily deep trees. No special API; just
\code{add} in a loop.

\subsection*{Composing with the DSL}

The output of a builder is an \code{Element}, so it composes with
the DSL:

\begin{lstlisting}
auto root = v(
    text(\"Header\") | Bold,
    render_tree(my_tree),       // a builder-produced Element
    text(\"Footer\") | Dim
);
\end{lstlisting}

You can switch between styles as the structure demands.

\section{Tradeoffs}

Builders are not free:

\begin{itemize}[leftmargin=*]
  \item \textbf{Allocation.} Each builder allocates the
    \code{vector<Element>} for its children eagerly. The DSL's
    compile-time path avoids this for fixed-size containers.
  \item \textbf{Type erasure.} The builder takes
    \code{Element} as the child type; the DSL's pipe overloads
    preserve more type information. If you need
    \code{constexpr}-correctness, the DSL is the path.
  \item \textbf{Error messages.} A typo in a setter
    (\code{borderr(\\ldots)}) is a normal C++ error: ``no member
    function''. The DSL's concept-based errors are sometimes
    clearer for structural mistakes.
\end{itemize}

The performance difference at frame time is usually negligible
(one vs zero allocations on a tree with hundreds of nodes is
microseconds). Pick by ergonomics, not by performance.

\begin{aside}
Internally, the DSL's \code{build()} produces the same
\code{Element} that a builder would. The DSL is a thin
compile-time wrapper that happens to skip some allocations; the
runtime model is the same.
\end{aside}

\section{Pitfalls}

\begin{warning}
\textbf{Forgetting to consume the builder.} If you do not call
\code{build()} (or pass it to \code{add()} / \code{widget()}),
the builder goes out of scope and its data is dropped.
\code{-Wunused-value} will warn if your builder is
\code{[[nodiscard]]}.
\end{warning}

\begin{warning}
\textbf{Reusing a moved-from builder.} After
\code{std::move(b).build()}, \code{b} is moved-from. Touching it
is technically allowed (the destructor runs) but reading it is
nonsense. The rvalue qualifier prevents this in practice.
\end{warning}

\begin{warning}
\textbf{Mixing units across setters.} \code{padding(1)} means
``all sides 1 cell''. \code{padding(1, 2)} means ``vert 1, horiz
2''. \code{padding(Padding\{.top=1, \\ldots\})} is the explicit
form. Read your code aloud; the unit should be obvious.
\end{warning}

\begin{warning}
\textbf{Reading layout from a half-built tree.} The builder does
not run layout; the rects are uninitialised until paint. If a
helper inspects \code{rect.w} on a builder-produced element
before painting, it will read zero.
\end{warning}

\section*{Workshop}

\begin{enumerate}[leftmargin=*]
  \item Write a function \code{Element render\_table(const
    std::vector<std::vector<std::string>>\& rows)} that uses
    builders to assemble a table.
  \item Read
    \file{maya/include/maya/element/builder.hpp} and identify
    every setter on \code{BoxBuilder}. Map each to a property
    on \code{FlexStyle}.
  \item Convert a small DSL expression into the equivalent
    builder code, then back. Note which form is more readable
    for each shape.
\end{enumerate}

\begin{recap}
Builders are the runtime spelling of the DSL. They produce the
same \code{Element} that the DSL produces; the difference is
where the work happens. Use builders when structure depends on
data; use the DSL when structure is fixed. Both compose at the
\code{Element} boundary.
\end{recap}
